\section{Data Collection}\label{sec:data_collection}

The orthographic and phonetic analyses in this paper are computed over a corpus
of Biblical translations that we collected ourselves. Because that corpus is
unpublished, its provenance cannot be established by citation, and we therefore
document the collection procedure here in full rather than referring the reader
elsewhere.

\subsection{Source and Procedure}

The corpus was drawn from YouVersion (\url{https://www.bible.com}), an online
Bible platform that hosts translations under license from their publishers. One
translation was selected per language for each of the sixteen languages
in~\cref{tab:ph_languages}. \Cref{tab:seeds} records the exact translation used
for each, together with the number of chapters obtained from it. Eight of the
translations are complete Bibles and were seeded at Genesis 1. The remaining
eight are New Testament translations still in progress and were seeded at
Matthew 1, which is why their chapter counts are uniformly lower.

\begin{table}[h!]
    \centering
    \footnotesize
    \begin{tabular}{llr}
        \hline
        \textbf{Language} & \textbf{Seed chapter}        & \textbf{Chapters} \\
        \hline
        Tagalog           & \texttt{2195/GEN.1.ABTAG01}  & 977               \\
        Cebuano           & \texttt{562/GEN.1.RCPV}      & 977               \\
        Ilocano           & \texttt{782/GEN.1.RIPV}      & 977               \\
        Hiligaynon        & \texttt{2190/GEN.1.MBBHIL12} & 1{,}104           \\
        Bikol             & \texttt{890/GEN.1.MBBBIK92}  & 1{,}104           \\
        Waray-Waray       & \texttt{2198/GEN.1.MBBSAM}   & 1{,}104           \\
        Kapampangan       & \texttt{1141/GEN.1.PMPV}     & 1{,}104           \\
        Pangasinan        & \texttt{1166/GEN.1.PNPV}     & 1{,}104           \\
        Adasen            & \texttt{2812/MAT.1.YBT}      & 209               \\
        Chavacano         & \texttt{1129/MAT.1.CBKNT}    & 209               \\
        Paranan           & \texttt{438/MAT.1.PRF}       & 209               \\
        Tausug            & \texttt{1319/MAT.1.TSG}      & 209               \\
        Romblomanon       & \texttt{2244/MAT.1.BKR}      & 209               \\
        Masbatenyo        & \texttt{1222/MAT.1.MSB}      & 209               \\
        Kinaray-a         & \texttt{1489/MAT.1.KRJNT}    & 209               \\
        Yami              & \texttt{2364/MAT.1.SNT}      & 209               \\
        \hline
                          & \textbf{Total}               & \textbf{10{,}123} \\
        \hline
    \end{tabular}
    \caption{Seed chapter and yield per language. Each seed is given relative to
        the prefix \url{https://www.bible.com/bible/}.}\label{tab:seeds}
\end{table}

YouVersion serves each chapter as an HTML document in which every span whose
class is prefixed with \texttt{ChapterContent\_verse\_\_} encloses one verse.
Within such a span, the verse text is carried by descendants prefixed
\texttt{ChapterContent\_content\_\_}, while editorial apparatus supplied by the
platform rather than by the translators, such as footnotes and
cross-references, occupies \texttt{ChapterContent\_note\_\_} descendants
interleaved with the text. We retained only the \texttt{content} descendants,
concatenated in document order, and discarded every sibling class. The corpus
consequently contains the translated text alone, and no editorial contribution
of YouVersion's own. Each chapter page also carries a navigation anchor to its
successor, which the collector followed from the seed, so each translation was
traversed as a linear chain rather than as a general crawl of the site.

Collection was performed once, in September 2025, using a purpose-written
crawler in Go built on the Colly framework. Traversal from each seed terminated
when the next-chapter link was exhausted rather than at the configured ceiling
of 30,000 chapters, which was never binding. Each language was collected in its
own goroutine, giving sixteen
concurrent request streams, and the implementation issued two requests per
chapter, one to extract the verses and a second to resolve the next-chapter
link. The sixteen streams together therefore placed on the order of 20,000
requests on the host. No crawl delay or rate limit was configured. The crawler
transmitted Colly's default user agent, which identifies the client as an
automated agent rather than imitating a browser.

Only pages that YouVersion serves publicly and anonymously were retrieved. The
crawler held no account, transmitted no credentials, and circumvented no
paywall, login, or other access control; it read the same markup that the site
returns to any unauthenticated visitor.

\section{Ethical Considerations}\label{sec:ethics}

Because the corpus underlying this paper was collected from a third party's
website rather than obtained from a published dataset, the procedure
of~\cref{sec:data_collection} requires scrutiny against the two instruments that
govern automated access to that site: its robots exclusion file and its Terms of
Use. We reviewed both, and we report the outcome in full, including where it is
unfavorable to us.

\subsection{Robots Exclusion}

At the time of review, the robots exclusion file at
\url{https://www.bible.com/robots.txt} applied a single rule set to all user
agents. It disallows \texttt{/bible/*/*/notes}, the user and plan search
endpoints, and five language-code directories, and it specifies no
\texttt{Crawl-delay}. The chapter pages we retrieved lie under \texttt{/bible/}
but match none of the disallowed patterns, and no disallowed path was requested.
The crawl therefore conformed to the robots exclusion file.

This conformance was, however, incidental rather than enforced. The Colly
collector sets its \texttt{IgnoreRobotsTxt} option to true on construction and
our implementation did not override it, so the crawler never retrieved or
evaluated the exclusion file at run time. Conformance was established by our own
later inspection, not by the software: had the file disallowed \texttt{/bible/},
the crawler would have proceeded regardless. We record this because a claim of
conformance that rests on coincidence is weaker than one that rests on
enforcement, and the reader is entitled to the distinction.

\subsection{Terms of Use}

The YouVersion Terms of Use are less permissive than the exclusion file. Under
the heading ``Permitted and Unpermitted Use'' they prohibit the user from

\begin{quote}
    ``us[ing] any robot, spider, or other automatic device, process, or means to
    access YouVersion for any purpose, including monitoring or copying any of the
    material on YouVersion''.
\end{quote}

We do not claim that our collection complies with this clause. It does not. The
crawler described in~\cref{sec:data_collection} is an automatic device that
accessed YouVersion for the purpose of copying material from it, which is the
conduct the clause names. A permissive robots exclusion file does not cure this,
as the two instruments are independent and the Terms are the governing agreement
between the site and its users. The order of events is likewise material and we
state it plainly: collection was carried out in September 2025, and the review
reported here was conducted afterwards. The Terms were not consulted before the
corpus was obtained, and the question was settled after the fact rather than in
advance.

We disclose this rather than assert a compliance we do not have. So that the
reader may weigh the matter independently, we set out the circumstances on both
sides.

Four circumstances mitigate. The use is non-commercial and academic, and no
revenue is derived from the corpus or from any work built upon it. No technical
protection measure was defeated: every page retrieved is served publicly and
anonymously, and the crawler bypassed no paywall, login, or access control. The
crawler identified itself honestly rather than imitating a browser, so its
automated character was apparent in the host's request log and was never
disguised. The collection was a single bounded pass rather than a standing
service, and nothing continues to poll the site.

Three circumstances aggravate, and we do not minimize them. No crawl delay was
configured, so requests were issued as fast as the parser allowed rather than at
a rate calibrated to the host's capacity. Sixteen concurrent streams ran
simultaneously. The implementation fetched each chapter twice, once for its
verses and once for its next-chapter link, which doubled our load on the host
for no benefit to us and could have been avoided by resolving both from a single
response. The approximately 20,000 requests we placed were thus roughly twice as
many as the task required, delivered without pacing. A defensible crawl would
have used one request per chapter, a fixed delay, and a single stream.

\subsection{Remediation}

The defect identified here is remediable, and the remedy is specific. YouVersion
operates a Platform API at \url{https://developers.youversion.com} which provides
programmatic access to its Bible collections under a developer agreement, at no
cost, and with a non-commercial restriction that an academic project satisfies by
construction. It covers substantially more versions and languages than the
sixteen used here. Obtaining the corpus through that interface would place the
collection within the terms the site sets for automated access, and would
dissolve the conflict described above rather than merely mitigate it. We
recommend that any continuation of this work re-collect through the Platform API,
and that the present corpus be treated as provisional until it does.
